\documentclass[10pt,letterpaper]{article}
\usepackage[T1]{fontenc}
\usepackage[utf8]{inputenc}
\usepackage{charter}
\usepackage{textcomp}
\usepackage[letterspace=60]{microtype}
\usepackage[margin=0.5in]{geometry}
\usepackage{enumitem}
\usepackage{xcolor}
\usepackage[hidelinks]{hyperref}
\definecolor{link}{HTML}{4A6FA5}
\pagestyle{empty}
\setlength{\parindent}{0pt}
\setlength{\parskip}{0pt}
\raggedbottom
\linespread{1.0}
\emergencystretch=2em

% ---------- section headings ----------
\newcommand{\cvsection}[1]{%
  \par\addvspace{7pt}%
  {\large\bfseries\MakeUppercase{#1}}\par\nobreak
  \vspace{-0.55\baselineskip}\rule{\linewidth}{0.6pt}\par\nobreak\vspace{1.5pt}}
\newcommand{\cvsub}[1]{%
  \par\addvspace{5pt}{\bfseries\textls{\MakeUppercase{#1}}}\par\nobreak\vspace{1.5pt}}

% ---------- org / role entries ----------
\newcommand{\entry}[4]{% org, location, role, dates
  \par\addvspace{2.5pt}%
  {\large\bfseries #1}\hfill #2\par\nobreak
  \textit{#3}\hfill #4\par\nobreak}
\newcommand{\mentor}[1]{\textit{#1}\par\nobreak}
\newenvironment{bullets}{%
  \begin{itemize}[leftmargin=0.95em,labelsep=0.35em,itemsep=0pt,parsep=0pt,topsep=1pt,partopsep=0pt,
    label={\raisebox{0.25ex}{\tiny$\bullet$}}]}{\end{itemize}}

% ---------- one-line items with year at right ----------
\newlength{\yw}\setlength{\yw}{0.78in}
\newcommand{\yitem}[2]{%
  \par\noindent\parbox[t]{\dimexpr\linewidth-\yw}{\raggedright\hangindent=1.6em\hangafter=1 #1\strut}%
  \hfill #2\par}
\newcommand{\pitem}[1]{\par\noindent{\raggedright\hangindent=1.6em\hangafter=1 #1\par}}

% ---------- publications ----------
\newenvironment{pubs}{%
  \begin{itemize}[leftmargin=2.6em,labelwidth=2.3em,labelsep=0.3em,align=left,itemsep=0pt,parsep=0pt,topsep=0pt]}{\end{itemize}}
\newcommand{\pub}[1]{\item[\textbf{[#1]}]}
\newcommand{\me}{\textbf{D. Pillis}}

\begin{document}

\begin{center}
  {\LARGE\bfseries\textls[80]{DANIEL PILLIS}}\\[3pt]
  {\color{link}\href{mailto:pillis@media.mit.edu}{pillis@media.mit.edu}
  \textcolor{black}{\,|\,} \href{https://dpillis.com}{dpillis.com}
  \textcolor{black}{\,|\,} \href{https://www.linkedin.com/in/pillis/}{linkedin.com/in/pillis}}
\end{center}
\vspace{-6pt}

\cvsection{Research Interests}
Hello! I am an interdisciplinary researcher and computer scientist. Innovative Technical Artist and Interaction
Designer with years of experience in software engineering, research in XR production, and interactive/immersive
media for virtual storytelling. A proven track record as a leader and collaborator in complex, cross-functional
environments.

\cvsection{Academic}
\entry{Massachusetts Institute of Technology}{Cambridge, MA}{Masters of Media Arts \& Sciences, MIT Media Lab}{2024}
\begin{bullets}\item Advisors: Pattie Maes, Hiroshi Ishii, Danielle Wood, Alan Lightman\end{bullets}
\entry{Carnegie Mellon University}{Pittsburgh, PA}{Masters of Fine Arts, focusing on Virtual Reality \& Robotics, School of Fine Arts}{}
\begin{bullets}\item Advisors: Jessica Hodgins, Chris Atkeson, Ivan Sutherland\end{bullets}
\entry{Rutgers University}{New Brunswick, NJ}{Bachelor of Arts, Cognitive Science, Visual Art, Literature}{}

\cvsection{Publications}
\cvsub{Journal Articles}
\begin{pubs}
\pub{J1} H. Ishii, \me*, P. Pataranutaporn*, X. Xiao, A. Algargoosh, L. Li, J.-B. Labrune. TeleAbsence: Visions of
  Afterlife Telepresence. \textit{Presence Journal (MIT Press)}, 2025.
\end{pubs}
\cvsub{Conference \& Workshop Papers}
\begin{pubs}
\pub{C5} \me, J. Ferguson. Return of the Jibo: Generative AI \& Social Robots for Virtual Production Education.
  \textit{MIT AI \& Education Conference}, Summer 2025.
\pub{C4} \me, P. Cerner, M. Sawtell. Real-Time Stable Diffusion for Virtual Production: Room-Scale Generative AI
  Towards the Ultimate Display. \textit{Generative AI for Design Workshop}, 2025. (Winner, Best Presentation)
\pub{C3} \me. Towards Temporal Telepresence \& Multisensory Spatial Intelligence: Reconstructing the Human
  Experience through TeleAbsence. \textit{IEEE VR, Workshop on Memory}, 2025.
\pub{C2} \me*, P. Pataranutaporn*, P. Maes, M. Sra. AI Comes Out of the Closet: Using AI-Generated Virtual
  Characters to Help Individuals Practice LGBTQIA+ Advocacy. \textit{ACM Conference on Intelligent User Interfaces
  (IUI)}, 2024.
\pub{C1} H. Ishii, O. Kilic Afsar, W. Babatain, C. Fang, J. Forman, S. Nicita, \me, L. Yang, J.-B. Labrune.
  Human--Material Interaction for Expression, Performance, and Remembering. \textit{CHI 2023 Workshop on
  Human--Material Interaction for Expression, Performance, and Remembering}, 2023.
\end{pubs}
\cvsub{Preprints \& Under Review}
\begin{pubs}
\pub{P2} C. Pierce, \me, P. Sinha, J. Skerswetat, S. Riskin, K. Zolot. Perceive More by Seeing Less: Immersive Vision
  Augmentation Systems to Support Vision Rehabilitation. \textit{Manuscript submitted for publication}, 2025.
\pub{P1} \me*, P. Pataranutaporn*, P. Maes, M. Sra. AI Theatre: Interactive Conversational Interfaces with Virtual
  Agents Can Increase Confidence and Self-Efficacy in Devising Solutions for Interpersonal Conflict Scenarios.
  \textit{Manuscript submitted for publication}, 2024.
\end{pubs}
\vspace{1pt}{\small\hfill * Equally contributing authors}

\cvsection{Academic Experience}
\entry{Emerson College}{Boston, MA}{Visiting Assistant Professor of Virtual Production}{2024 -- Present}
\begin{bullets}
\item Pioneering research in next-generation virtual production, leveraging a state-of-the-art 2240 $\times$ 1440, 1.9mm-pitch LED volume to push the boundaries of real-time, cinematic immersive media
\item Lead developer of large-scale, interactive immersive environments across Boston and Los Angeles campuses
\item Publishing research at the intersection of generative media, immersive environments, and human/AI interaction, advancing new paradigms for AI-driven storytelling and spatial computing
\item Mentor and instructor for graduate and undergraduate students, establishing industry-leading best practices in virtual production, AI-assisted filmmaking, and real-time workflows
\item Building high-fidelity, experimental interactive systems and research prototypes in Unreal Engine and Unity, focused on real-time rendering, immersive interaction, and advanced virtual production pipelines
\item Classes taught: AI Filmmaking, Topics in Virtual Production
\end{bullets}

\entry{Massachusetts Institute of Technology, MIT Media Lab}{Cambridge, MA}{Research Associate}{2022 -- 2024}
\begin{bullets}
\item Designed, fabricated, and programmed advanced pneumatic and mixed-material tangible interfaces, combining custom PCBs, soft-robotic actuation, and embedded microcontrollers (Arduino, FlowIO), evaluated through human-subject user studies
\item Engineered AR/VR simulations, digital-twin pipelines, and Gaussian splatting / photogrammetry workflows for high-fidelity interactive spatial media, enabling novel memory-support and autobiographical interfaces
\item Published peer-reviewed research in top-tier international venues, including ACM IUI, ACM CHI, and IEEE VR, with a second-author journal publication in MIT Presence Journal
\item Co-led the TeleAbsence research initiative under the supervision of Hiroshi Ishii, overseeing \$40K+ in research funding and delivering dozens of successful industry-facing demos and installations
\item Captured and processed hundreds of personal objects and heirlooms via high-resolution 3D scanning, constructing a large-scale multimodal dataset to support memory-driven and autobiographical AI model development
\item Conceived and delivered public-facing outreach exhibitions, XR installations, and collaborative research projects spanning Tangible Media Group, Fluid Interfaces, City Science, and DesignX
\item Classes taught: MAS.834: Tangible User Interfaces with Dr. Hiroshi Ishii; MAS.S61: Metaverse with Dr. Joseph Paradiso, Hiroshi Ishii \& Valentin Heun; MAS.S67: Design for the Future with Dr. Danielle Wood
\end{bullets}

\entry{Academy Software Foundation}{}{Graduate Intern}{2023}
\mentor{Supervisor: Wan Chun Alex Ma}
\begin{bullets}
\item Led a research project overseen by researchers from Industrial Light \& Magic and Meta, showcasing emerging applications of generative AI in virtual production
\item Built a browser-based, motion-capture--driven simulation using JavaScript, Three.js, and WebSockets, integrating GPT-3--based narrative generation with parameterizable real-time mocap behaviors
\item Published peer-reviewed outcomes at the ACM Intelligent User Interfaces Conference
\end{bullets}

\entry{MIT DesignX}{Cambridge, MA}{Startup Founder, SoftDesk}{2022}
\begin{bullets}
\item Founded and led SoftDesk, a sensory-computing startup developing tangible, multisensory interfaces through soft robotics, pneumatic actuation, and embedded sensing
\item Secured \$10K in competitive seed funding and led a four-person interdisciplinary team from Harvard GSD, driving product--market fit, rapid prototyping sprints, and milestone-based execution
\item Delivered investor-facing demos and pitch decks to audiences of hundreds of venture capitalists, clearly articulating technical architecture, IP strategy, and market opportunities
\end{bullets}

\entry{Princeton University}{Princeton, NJ}{Associate Professional Specialist / Visiting Instructor}{2020 -- 2021}
\begin{bullets}
\item Developed VR-centric curricula for remote art and technology courses, covering creative coding, 3D animation, robotics programming, and VR/AR/MR pipelines (Unity, C\#, shaders)
\item Built a networked Unity simulation of the Princeton Art Museum for Zoom-based instruction, enabling remote students to co-inhabit and annotate shared 3D spaces in real time
\item Prepared and managed course technology infrastructure, including custom Immersive Media Kits and Unity-ready mini-PCs for advanced real-time development
\item Classes taught: THR STC 210: Storytelling with Technology; STC 209: Engineering Transformations
\end{bullets}

\entry{Virginia Tech}{Blacksburg, VA}{Research Assistant Professor in Immersive Environments}{2018 -- 2020}
\begin{bullets}
\item Grant-funded researcher and creative technologist, leading cross-disciplinary collaborations across art, engineering, and the sciences
\item Project lead for ``The Cube,'' a \$15M, one-of-a-kind immersive VR research facility, featuring 150 spatial audio channels, 32 Qualisys motion-capture cameras, and 360\textdegree{} projection environments
\item VR/AR hardware specialist and technical consultant, supporting Oculus/Quest, HTC Vive, and HoloLens, alongside immersive video and projection-mapping workflows
\item Designed and taught Honors-level Unity courses, integrating animation, motion capture, procedural generation, and immersive storytelling for both art and STEM contexts
\item Classes taught: UH 3004: Immersive Virtual Environments for Art, Data, and Research
\end{bullets}

\entry{ACM SIGGRAPH International Conference}{}{SIGGRAPH Diversity \& Inclusivity Committee}{2018 -- 2025}
\begin{bullets}
\item Collaborated with industry peers to advance diversity, equity, and inclusion across computer graphics and interactive technology communities
\item Planned and led panels, presentations, and conference programming addressing gender theory, critical race theory, queer studies, intersectionality, disability studies, and related frameworks
\item Hosted and contributed to annual Diversity \& Inclusion Summits, fostering actionable dialogue and best practices for inclusive cultures in XR, VR, animation, interactive graphics, and CGI industries
\end{bullets}

\entry{Gujarat Science Center}{Ahmedabad, India}{Robotics Education Consultant}{2018}

\entry{University of Maine, Innovation Media Research Center}{Orono, ME}{Researcher-in-Residence}{2018}
\begin{bullets}
\item Demonstrated interactive robotics and AI systems, including SoftBank NAO--based installations, in an Artificial Intelligence and Art course co-taught with Sofian Audry
\item Developed ARKit-based iOS prototypes in Unity and delivered hands-on instruction in augmented reality development
\item Advised graduate students and consulted on the design and development of a Virtual Reality Lab, supporting research and curriculum planning
\end{bullets}

\entry{Carnegie Mellon University, School of Computer Science}{Pittsburgh, PA}{Research Associate, The Robotics Institute}{2016 -- 2018}
\begin{bullets}
\item Visiting Researcher in Christopher Atkeson's soft robotics laboratory
\item Developed interactive robotics--VR systems using historic robot primitives, enabling users to design and simulate custom robotic behaviors in virtual reality
\item Cataloged, digitized, and contextualized a historic robotics collection, designing a taxonomy and framework to generate a structured 3D dataset
\item Produced high-quality instructional videos and animations on robotics concepts and programming for K--12 and university outreach, broadening participation in robotics and computer science
\item Accepted proposal to present an overview of this work at IROS Robotics Conference, Vancouver, BC
\end{bullets}

\entry{Carnegie Mellon University, Posner Center Rare Book Collection}{Pittsburgh, PA}{Curatorial Fellow}{2015 -- 2016}
\begin{bullets}
\item Collaborated with pioneering computer scientist Ivan Sutherland to document the invention of the first virtual reality headset, designing and implementing an interactive VR application
\item Conducted interdisciplinary research spanning computer science, computer graphics, and automata, bridging theory and interactive systems
\item Curated a large-scale public exhibit integrating rare books, robotics artifacts, and an embedded interactive VR experience; oversaw visual design, printed materials, and exhibition documentation
\item Coordinated invited lectures by Marc Raibert (Boston Dynamics) and Ivan Sutherland, jointly hosted by the School of Art and the Robotics Institute
\end{bullets}

\entry{Carnegie Mellon University, School of Art / School of Computer Science}{Pittsburgh, PA}{Graduate Research Assistant}{2013 -- 2016}
\begin{bullets}
\item Developed and deployed emerging-technology workflows to support faculty in teaching animation and advanced digital production
\item Instructed students in motion-capture pipelines, including capture, rigging, and retargeting of performance data
\item Taught core animation and production workflows, spanning storyboarding, video editing, 3D modeling, and animation in Maya, Blender, and ZBrush
\item Supported student mastery of advanced techniques, including cloth simulation, render layering, character rigging, and physically based smoke, fire, and dynamics simulations
\end{bullets}

\cvsection{Professional Experience}
\yitem{\textbf{Federico Solmi Studio}, Virtual Reality Developer, New York, NY}{2017 -- 2018}
\yitem{\textbf{Artsy}, Virtual Reality Consultant, New York, NY}{2016}
\yitem{\textbf{The Wooster Group}, Theatre Assistant Technician, Performing Garage, New York, NY}{2016}
\yitem{\textbf{Skowhegan Residency}, Artist's Assistant, Maine}{2014}
\yitem{\textbf{Elizabeth Dee Gallery}, Digital Archivist, New York, NY}{2011 -- 2013}

\cvsection{Teaching}
\yitem{Virtual Production, \textbf{Emerson College}}{2024 -- 2025}
\yitem{AI Filmmaking, \textbf{Emerson College}}{2024 -- 2025}
\yitem{Tangible Interfaces, Fall 2023, Graduate Assistant, \textbf{Massachusetts Institute of Technology, Media Lab}, Cambridge, MA}{2023}
\yitem{MAS.S61: The Metaverse: What, How, Why, and When, Spring 2023, Graduate Assistant, \textbf{Massachusetts Institute of Technology, Media Lab}, Cambridge, MA}{2023}
\yitem{Transformations in Engineering \& the Arts, Spring 2021, \textbf{Princeton University}, Princeton, NJ}{2021}
\yitem{Technology \& Storytelling for Performance, Teaching Assistant, Fall 2020, \textbf{Princeton University}, Princeton, NJ}{2020}
\yitem{UH 3004: Immersive Virtual Environments for Art, Data, and Research, \textbf{Virginia Tech, Honors College}, Blacksburg, VA}{2019}
\yitem{Animation, Art \& Technology, Spring 2016, Teaching Assistant, Dr. Jessica Hodgins \& James Duesing, \textbf{Carnegie Mellon University}, Pittsburgh, PA}{2016}
\yitem{Technical Character Animation, Fall 2015, Teaching Assistant, Assistant Prof. Spencer Diaz, \textbf{Carnegie Mellon University}, Pittsburgh, PA}{2015}
\yitem{Experimental Animation, Spring 2014, Teaching Assistant, Professor Jessica Hodgins (C.S.), Professor James Duesing (Art), \textbf{Carnegie Mellon University}, Pittsburgh, PA}{2014}
\yitem{Foundations Concept Studio; Space and Time, Fall 2014, Teaching Assistant, Visiting Prof. Jonathan Armistead, \textbf{Carnegie Mellon University}, Pittsburgh, PA}{2014}
\yitem{Electronic Media Studio: Animation \& Video, Spring 2014, Teaching Assistant, Prof. Paolo Pederecini, \textbf{Carnegie Mellon University}, Pittsburgh, PA}{2013}

\cvsection{Presentations and Exhibitions}
\cvsub{Talks}
\yitem{Invited Speaker, AI Comes out of the Closet, \textbf{Hubspot}}{2025}
\yitem{Invited Speaker, \textbf{Mark Billingshurst}}{2025}
\yitem{Invited Speaker, \textbf{IEEE VR}, St Malo, France}{2025}
\yitem{Invited Speaker, \textbf{HarvardXR}}{2025}
\yitem{Invited Speaker, Panel with Alvin Wang Graylin, \textbf{MIT AI Film Hackathon}}{2025}
\yitem{Invited Speaker, \textbf{SIGGRAPH Sparks Talk}}{2025}
\yitem{Invited Speaker, \textbf{SIGGRAPH Village Talks}}{2024}
\yitem{Invited Speaker, \textbf{MIT Media Lab}, Dr. Danielle Wood}{2024}
\yitem{Invited Speaker, \textbf{SIGGRAPH International Conference}, Los Angeles, CA}{2023}
\yitem{Invited Speaker, \textbf{AR/VR Hackathon, Harvard University}, Cambridge, MA}{2023}
\yitem{Conference on Adaptive Technology, \textbf{Digital IDEAS Summer Institute, University of Michigan}, Ann Arbor, MI}{2022}
\yitem{Invited Speaker, DEI Panel Discussion, \textbf{High-Performance Graphics Conference}}{2021}
\yitem{Diversity, Equity \& Inclusion Volunteer, \textbf{SIGGRAPH 2021}}{2021}
\yitem{SIGGRAPH LGBTQ+ Meetup, Augmented Reality Drag Show using AR via Zoom, \textbf{ACM SIGGRAPH International Conference}}{2021}
\yitem{Invited Speaker, IDEA: Inclusion, Diversity, Equity, and Accessibility, \textbf{Princeton ResInDe Virtual Conference, Princeton University}}{2021}
\yitem{Animation collaboration with musician Bora Yoon, \textbf{Living at the Intersection Symposium, Princeton University}}{2020}
\yitem{Visiting Lecture, Summer Course, \textbf{UCLA Digital Media Arts} (invited by Miles Peyton)}{2019}
\yitem{Visiting Lecture, Imitation of Life, Virtual Reality Workshop, \textbf{Cooper Union}, New York, NY (invited by Joao Enxuto)}{2019}
\yitem{Invited Lecture, Gender Representation in Computer Graphics, \textbf{SIGGRAPH Asia}, Brisbane, Australia}{2019}
\yitem{Panel member \& organizer, Rendering Gender, with Jacob Gaboury \& Bo Ruberg, \textbf{SIGGRAPH Conference}, Los Angeles, CA}{2019}
\yitem{Organizer, Diversity and Inclusion Summit, \textbf{SIGGRAPH Conference}, Vancouver, BC}{2018}
\yitem{Invited Lecture, Robotics and the History of the Humanoid, \textbf{University of Maine IMRC}, Orono, MN}{2018}
\yitem{Invited Lecture, Introduction to Augmented Reality using ARKit, Workshop, \textbf{University of Maine IMRC}, Orono, MN}{2018}
\yitem{Visiting Lecturer, Artificial Intelligence and Art, with Prof. Sofian Audry, \textbf{University of Maine IMRC}, Orono, MN}{2018}
\yitem{Invited Lecture, Towards an Artist in the Lab Framework, Workshop, \textbf{iRos Robotics Conference}, Vancouver, BC}{2017}
\yitem{Science Educator, with Dr. Chris Atkeson, \textbf{Robotics Institute, Carnegie Mellon University}, Pittsburgh, PA}{2017}
\yitem{Visiting Lecturer, Figure Drawing for Humanoid Robots, \textbf{Columbia University School of the Arts}, New York, NY}{2017}
\yitem{Visiting Critic, Physical Computing, \textbf{Carnegie Mellon University}, Pittsburgh, PA}{2017}
\yitem{Visiting Critic, Senior Studio, \textbf{Carnegie Mellon University}, Pittsburgh, PA}{2016}
\yitem{Visiting Lecturer, Six X Eight Artist Lecture Series, \textbf{Silver Eye Center for Photography}, Pittsburgh, PA}{2016}
\yitem{Committee Member, Public Arts Project Selection, \textbf{Open Engagement}, Pittsburgh, PA}{2015}
\yitem{Undergraduate Review Committee, \textbf{Carnegie Mellon University}}{2015}

\cvsub{Exhibitions}
\yitem{Three Body Problem, curated by Ashik Zaman and Daniel Pillis, \textbf{Harvard Conflux, Stora Galleriet}, March 20 -- April 12}{2025}
\yitem{MIT After Dark, Vision Simulations, with Cayden Pierce, \textbf{MIT Museum}}{2025}
\yitem{Design for the Future, \textbf{MIT Media Lab}, Massachusetts Institute of Technology, Cambridge, MA}{2023}
\yitem{Media Lab Gender Bender, \textbf{MIT Media Lab}, Massachusetts Institute of Technology, Cambridge, MA}{2023}
\yitem{Virtual Princeton Art Museum Exhibit, online virtual experience of \textbf{Princeton University Art Collections}}{2020}
\yitem{Federico Solmi: The Bacchanalian Ones, \textbf{Rowan University}, NJ (Freelance VR Developer)}{2020}
\yitem{Boo Box: VR installation, \textbf{New Orleans Film Festival, Cinema Reset VR Exhibit}, curated by Rachel Lin Weaver, New Orleans, LA}{2019}
\yitem{New \& Improved, Group Exhibit curated by Sophie Sobers, \textbf{Index Art Center}, Newark, NJ}{2019}
\yitem{Dancing Plants, Science Festival, \textbf{The Institute for Creativity, Art, and Technology, Virginia Tech}, Blacksburg, VA}{2019}
\yitem{Visual Language of Chromatin Architecture, ICAT Day, \textbf{The Institute for Creativity, Art, and Technology, Virginia Tech}, Blacksburg, VA}{2019}
\yitem{Curator, Open at the Source, Exhibit on the intersection between Art \& Science, \textbf{The Institute for Creativity, Art, and Technology, Virginia Tech}, Blacksburg, VA}{2019}
\yitem{History of Computer Graphics, Interactive Animation, \textbf{SIGGRAPH Art Gallery}, Vancouver, Canada}{2018}
\yitem{The Blue Plate, \textbf{University of Mary Washington}, Fredericksburg, VA}{2018}
\yitem{Mainframe, \textbf{Mint Museum}, Charlotte, NC}{2018}
\yitem{Robot Museum, \textbf{Interstate Projects}, Brooklyn, NY}{2017}
\yitem{Et Tu, Art Brute?, \textbf{Andrew Edlin Gallery}, New York, NY}{2017}
\yitem{Grandmother's House, \textbf{Space Gallery}, Pittsburgh, PA}{2016}
\yitem{Virtual Newell/Simon Simulation, \textbf{Miller Gallery, Carnegie Mellon MFA Thesis Show}, Pittsburgh, PA}{2016}
\yitem{Ivan Sutherland's Trojan Cockroach, \textbf{Posner Center, Carnegie Mellon}, Pittsburgh, PA}{2016}
\yitem{Grandmother's House, \textbf{Little Berlin}, Philadelphia, PA}{2015}
\yitem{An Evening with Cara Benedetto, Performance, \textbf{Museum of Contemporary Art}, Cleveland, OH}{2015}
\yitem{The Blue Plate: History of Computer Graphics, \textbf{Carnegie Mellon University}, Pittsburgh, PA}{2015}
\yitem{Intergalactic Immigration Office, curated by Jen Delos Reyes, \textbf{Open Engagement 2015}, Pittsburgh, PA}{2015}
\yitem{Andy Warhola's Living Room, \textbf{The Warhol Museum}, Pittsburgh, PA}{2014}
\yitem{City of Lost Men, \textbf{Kresge Theatre, Carnegie Mellon University}, Pittsburgh, PA}{2014}
\yitem{Grandmothers House, \textbf{4022 Woolslayer Way}, Pittsburgh, PA}{2013}
\yitem{Digital Archive 2011--2013, curated by Rebecca Jampol, \textbf{Newark Penn Station}, Newark, NJ}{2013}
\yitem{Queer Profiles in Courage, curated by Christopher Mitchell, \textbf{Leslie/Lohman Museum of Gay \& Lesbian Art}, NY, NY}{2012}
\yitem{Mix23 NYC Experimental LGBTQ Film Festival, \textbf{Theatre for the New City}, NY, NY}{2010}
\yitem{Last Days of the Iron Maiden, \textbf{Gallery Aferro}, Newark, NJ}{2009}
\yitem{Mason Gross BFA, \textbf{Rutgers University}, New Brunswick, NJ}{2009}
\yitem{One City, \textbf{Jajo Gallery}, Newark, NJ}{2009}
\yitem{Multiformity: The Art of Disease, curated by Deana Haggag, \textbf{Jajo Gallery}, Newark, NJ}{2009}

\cvsub{Residencies}
\yitem{\textbf{Media Archaeology Lab}, the University of Colorado at Boulder (forthcoming)}{2021}
\yitem{\textbf{University of Maine, Innovative Media, Research and Commercialization Center}, Researcher in Residence, Orono, ME}{2018}
\yitem{\textbf{Interstate Projects}, Brooklyn, NYC, NY}{2017}
\yitem{\textbf{The Robotics Institute, Carnegie Mellon University}, Pittsburgh, PA}{2016}

\cvsection{Technical Skills}
\pitem{\textbf{Software:} Unity, Unreal, Blender, Maya, ZBrush, Autodesk Software, Adobe Suite, Azure Kinect}
\pitem{\textbf{Hardware:} Oculus Quest 1--3, HoloLens, Magic Leap, Leap Motion, Arduino, Raspberry Pi, EEG, Computer Vision}
\pitem{\textbf{General Skills:} C\#, Python, Ubuntu, Windows 11, Mac OS, Networking, 3D Printing, HTC Vive, Microcomputers, Motion Capture, Project Planning, Asset Management}

\cvsection{Honors and Awards}
\yitem{\textbf{Best Presentation}, Generative AI for Design Workshop, MIT}{2025}
\yitem{\textbf{MIT Seed Fund}, \textbf{\$40K} award for TeleAbsence}{2024}
\yitem{\textbf{ICAT Major SEAD Grant}, Dancing Plants, Virginia Tech}{2018}
\yitem{\textbf{ICAT SEAD Grant}, Visual Language of Chromatin Architecture, Virginia Tech}{2018}
\yitem{Steelcase Education Grant, Virginia Tech}{2018}
\yitem{Innovation with Impact, Graduate Research Grant, Carnegie Mellon University}{2016}
\yitem{Rembacher Chamber Music Award, Carnegie Mellon University}{2016}
\yitem{Studio for Creative Inquiry Grant, Carnegie Mellon University}{2015}
\yitem{Full Tuition Merit Based Scholarship, Carnegie Mellon University}{2015}
\yitem{Public Art Scholarship, Rutgers University}{2014}

\cvsection{Select Media Coverage}
\yitem{Berkeley Beacon, ``How Emerson is---and isn't---using artificial intelligence in the classroom,'' by Y. Toby}{2025}
\yitem{MIT News, ```AI comes out of the closet' seeks to merge AI and LGBTQIA+ advocacy in a way as innovative as it is necessary,'' by David Sweeney}{2024}
\yitem{MIT News, ``Communicating Across Time,'' by Ken Schulman}{2022}
\yitem{Pittsburgh City Paper, ``CMU launches Project showcasing the History of Robotics at the University,'' by Lauryn Nania}{2021}
\yitem{NEXTPittsburgh, ``Robotics Project traces the roots of today's robots to CMU,'' by Michael Machosky}{2021}
\yitem{Ground Truth, ``Carnegie Mellon University Launches The Robotics Project to Preserve and Promote the Field's Legacy,'' by GT Staff}{2021}
\yitem{TribLive, ``Carnegie Mellon's Robotics Project aims to preserve History of Robotics with large-scale Archive,'' by Julia Felton}{2021}
\yitem{Pittsburgh Post Gazette, ``From trading cards to `cannibilized' robotic parts, CMU embarks on a mission to archive the Robotics Field,'' by Lauren Rosenblatt}{2021}
\yitem{Virginia Tech News, ``The Sound of Plants Dancing,'' by Max Esterhuizen}{2020}
\yitem{The Virginian Pilot, ``Virginia researchers are studying `the sound of plants dancing' to better the future of agriculture,'' by Katherine Hafner}{2020}
\yitem{Virginia Tech, Winter Course Feature}{2020}
\yitem{Vice, Creators Project, ``An Artist built an exact replica of Artificial Intelligence's Birthplace,'' by Sean Neumann}{2017}
\yitem{Digital Trends, ``Artist re-creates the birthplace of Artificial Intelligence,'' by Dyllan Furness}{2017}
\yitem{Atlas Obscura, ``Where Should Robots Go When They Retire?,'' by Sarah Laskow}{2016}
\yitem{iProgrammer, ``Sutherland's Trojan Cockroach On Show,'' by David Conrad}{2016}
\yitem{Pittsburgh City Paper, ``The Year in Visual Art,'' by Robert Raczka}{2015}
\yitem{The Warhol Blog, ``Andy Warhola's Living Room''}{2014}
\yitem{Pittsburgh Post Gazette, ``Grandma's belongings live as art installation in Lawrenceville,'' by Donna Nelson-Jones}{2014}

\cvsection{Service}
\cvsub{Mentoring}
\yitem{\textbf{Harvard Graduate School of Design Conflux Residency Program}, Mentor}{2024 -- Present}
\yitem{\textbf{Out in Tech}, Mentor}{2024}
\yitem{\textbf{MIT Reality Hack}, Mentor}{2023}
\yitem{\textbf{MIT MAS SOS}, mentored a total of 4 students from underrepresented groups in preparing their graduate applications to MIT}{Fall 2023}
\yitem{\textbf{Out in Tech Mentorship}, mentored 2 undergraduate students in 2020, 1 in 2021, and 1 in 2022, including discussing preparation for graduate school and the tradeoffs between academic and industry careers}{2022 -- 2024}
\cvsub{Advising}
\pitem{Minyu Chen, Emerson College Thesis Committee}
\pitem{Nikita Shokhov, Mixed Reality Independent Study, with Aisling Kelliher, Masters in Creative Technologies, School of Visual Arts, Virginia Tech}
\pitem{Mark Manual, Master's Thesis, Department of Computer Science, Virginia Tech}
\pitem{Yao Xiao, Advisor, Department of Computer Science, Master of Science, Virginia Tech}
\cvsub{Professional Affiliations}
\pitem{Special Interest Group on Computer Graphics and Interactive Techniques, SIGGRAPH}
\pitem{Association for Computing Machinery}

\end{document}
